% ---------------------------------------------------------------------------
\section{Fundamentação Teórica}
\label{sec:fundamentacao}
% ---------------------------------------------------------------------------

Este capítulo apresenta os fundamentos teóricos que sustentam o estudo da geração automática de fases tridimensionais para jogos digitais. A exposição organiza-se de modo progressivo: parte-se da definição de fases, de \textit{Level Design} e da distinção conceitual entre organização lógica e organização espacial; em seguida, abordam-se os princípios da Geração Procedural de Conteúdo (PCG), sua vertente baseada em aprendizado de máquina (PCGML) e a caracterização dos Modelos de Linguagem de Grande Porte (LLMs), incluindo suas limitações de raciocínio e ancoragem espacial. Por fim, discutem-se a modelagem lógica de fases por meio da Teoria dos Grafos, a representação modular e geométrica de \textit{layouts} tridimensionais, as arquiteturas híbridas guiadas por restrições e os critérios formais de validade, navegabilidade e solvabilidade.

% ---------------------------------------------------------------------------
\subsection{Fases de Jogos, \textit{Level Design} e \textit{Layouts} 3D}
\label{sec:level_design}
% ---------------------------------------------------------------------------

\subsubsection{Fases, Níveis e \textit{Level Design}}
\label{sec:fases_niveis}

No contexto dos jogos digitais, uma fase --- também denominada nível (\textit{level}) ou mapa --- constitui a unidade espacial, temporal e interativa na qual as regras e mecânicas do sistema ludológico se manifestam para o jogador \cite{tanagra, pcgbook}. Mais do que um cenário estático, uma fase opera como um contêiner estruturado para a experiência de jogo (\textit{gameplay}), delimitando o território de exploração, os recursos disponíveis, os obstáculos a serem transpostos e as condições de vitória ou transição \cite{tanagra, linden2014dungeons}. Em gêneros como ação e aventura, jogos de plataforma e \textit{Role-Playing Games} (RPGs), a qualidade da experiência interativa depende diretamente de como esses espaços são construídos e articulados \cite{dormans2010, linden2014dungeons}.

A disciplina responsável pelo planejamento, concepção e estruturação de fases é o \textit{Level Design} \cite{dormans2010, tanagra}. Seus objetivos fundamentais abrangem a condução da atenção e do deslocamento do jogador, a introdução gradual de mecânicas, o balanceamento entre desafio e recompensa e a regulação do ritmo da experiência (\textit{pacing}) \cite{dormans2010, inproceedings, linden2014dungeons}. Em fases do tipo masmorra (\textit{dungeon}), por exemplo, o \textit{Level Design} busca conciliar a liberdade de exploração espacial com um controle rigoroso sobre a progressão de desafios e quebra-cabeças (\textit{puzzles}), assegurando que o jogador não encontre obstáculos intransponíveis para o conjunto de habilidades ou recursos que detém em determinado instante da partida \cite{linden2014dungeons, shaker2016constructive}.

\subsubsection{\textit{Layout} de uma Fase}
\label{sec:layout_fase}

O \textit{layout} de uma fase refere-se ao arranjo estrutural e espacial dos elementos que compõem o ambiente jogável \cite{dormans2010, linden2014dungeons}. Em ambientes internos e masmorras, o \textit{layout} é classicamente decomposto em unidades arquitetônicas discretas, destacando-se as salas (ou câmaras), os corredores (ou vias de transição), os conectores (portas, portais, escadas e elevadores) e as regiões funcionais \cite{linden2014dungeons, shaker2016constructive, xu2025databasedrivenframework3dlevel}. As salas atuam tipicamente como arenas de interação concentrada --- abrigando combates, enigmas, recursos ou eventos narrativos ---, ao passo que os corredores estabelecem as ligações físicas entre essas áreas, modulando o tempo de travessia e a visibilidade entre setores adjacentes \cite{inproceedings, shaker2016constructive}.

Além da geometria estrutural que delimita o espaço navegável (pisos, paredes, tetos e desníveis), o \textit{layout} compreende a disposição interna de elementos funcionais e estéticos, frequentemente denominados instalações ou objetos de cena (\textit{facilities} ou \textit{props}), bem como entidades dinâmicas, tais como personagens não jogáveis (NPCs), inimigos, armadilhas e itens coletáveis \cite{linden2014dungeons, xu2025constraint}. Em ambientes tridimensionais (3D), a especificação de um \textit{layout} envolve não apenas a distribuição planar em dois eixos, mas também a verticalidade, a sobreposição de múltiplos pavimentos, as relações de linha de visão (\textit{line of sight}) e as dimensões volumétricas exatas de cada componente no espaço contínuo \cite{xu2025constraint, xu2025databasedrivenframework3dlevel}.

\subsubsection{Organização Lógica e Organização Espacial}
\label{sec:organizacao_logica_espacial}

Um princípio teórico central na literatura de \textit{Level Design} consiste na distinção analítica entre duas estruturas complementares que coexistem em uma fase: a \textit{missão} (organização lógica) e o \textit{espaço} (organização espacial) \cite{dormans2010, linden2014dungeons}. A missão corresponde à sequência lógica ou parcialmente ordenada de tarefas, ações e objetivos que o jogador deve cumprir para concluir o nível, tais como derrotar um guardião, obter uma chave específica, desativar um bloqueio e alcançar a saída \cite{dormans2010, inproceedings}. O espaço, por sua vez, refere-se à realização física e geométrica do ambiente virtual onde essas tarefas estão situadas, englobando as dimensões das salas, as distâncias métricas e a configuração arquitetônica das conexões \cite{dormans2010, linden2014dungeons}.

Embora o espaço físico acomode a missão e esta seja mapeada sobre aquele, as duas estruturas são conceitualmente independentes e regidas por princípios distintos \cite{dormans2010}. Uma missão estritamente linear --- na qual os objetivos precisam ser cumpridos em uma ordem única e fixa --- pode ser distribuída sobre um espaço físico não linear, como um arranjo em estrela (\textit{hub-and-spoke}) ou em anel, exigindo que o jogador retorne a áreas previamente visitadas à medida que desbloqueia novos acessos \cite{dormans2010}. Inversamente, uma missão não linear, dotada de múltiplos desafios paralelos ou caminhos alternativos, pode ser acomodada em uma geometria espacial simples \cite{dormans2010, inproceedings}. Essa dissociação teórica demonstra que a coerência da progressão de um nível reside primariamente em sua estrutura lógica de dependências, enquanto sua viabilidade física e navegacional depende da correta realização geométrica no espaço tridimensional \cite{dormans2010, inproceedings, linden2014dungeons}.

% ---------------------------------------------------------------------------
\subsection{Geração Procedural de Conteúdo}
\label{sec:pcg}
% ---------------------------------------------------------------------------

\subsubsection{Conceito e Objetivos da PCG}
\label{sec:conceito_pcg}

A Geração Procedural de Conteúdo (PCG --- \textit{Procedural Content Generation}) é definida na literatura como a criação algorítmica de conteúdos para jogos digitais de maneira autônoma ou com intervenção humana limitada e indireta \cite{pcgbook, summerville2018procedural, maleki2024procedural}. O escopo do termo ``conteúdo'' abrange praticamente todos os artefatos que integram um jogo, incluindo níveis, mapas, terrenos, malhas vegetais, texturas, efeitos sonoros, narrativas, missões (\textit{quests}) e conjuntos de regras, distinguindo-se do comportamento em tempo de execução dos agentes de inteligência artificial que controlam NPCs \cite{pcgbook, summerville2018procedural, maleki2024procedural}.

Historicamente introduzida em títulos pioneiros da década de 1980, como \textit{Rogue} (1980) e \textit{Elite} (1984), para contornar severas restrições de memória e armazenamento por meio da síntese algorítmica em tempo de execução \cite{pcgbook, summerville2018procedural, araujo2018geracao}, a PCG consolidou-se tanto na indústria quanto na pesquisa acadêmica em função de múltiplos objetivos \cite{linden2014dungeons, pcgbook}. Dentre as principais motivações para o seu emprego destacam-se: (i) a automação e a redução do tempo e do custo financeiro associados à autoria manual de ativos em larga escala; (ii) o aumento da variabilidade e da diversidade dos ambientes gerados, ampliando o fator de rejogabilidade (\textit{replayability}); (iii) o suporte à adaptabilidade centrada no jogador, permitindo ajustar dinamicamente os desafios às habilidades observadas; e (iv) a exploração de novas estéticas e mecânicas emergentes \cite{linden2014dungeons, pcgbook, summerville2018procedural}.

A avaliação e o projeto de geradores procedurais apoiam-se em propriedades técnicas fundamentais, notadamente a reprodutibilidade, a expressividade e a controlabilidade \cite{tanagra, linden2014dungeons, maleki2024procedural}. A reprodutibilidade decorre frequentemente do uso de Geradores de Números Pseudoaleatórios (PRNGs), nos quais a especificação de um valor inicial --- denominado semente (\textit{seed}) --- permite reconstruir de maneira determinística a mesma sequência de decisões algorítmicas \cite{maleki2024procedural}. A expressividade, frequentemente analisada por meio do espaço expressivo (\textit{expressive range}), caracteriza a variedade e as propriedades estilísticas ou estruturais dos artefatos que o gerador é capaz de produzir \cite{tanagra, inproceedings}. A controlabilidade, por fim, refere-se ao grau de domínio oferecido ao projetista ou ao sistema para direcionar intencionalmente as características do conteúdo gerado por meio de parâmetros, regras ou restrições, garantindo resultados consistentes sem sacrificar a diversidade \cite{linden2014dungeons, shaker2016constructive}.

\subsubsection{Principais Abordagens de PCG}
\label{sec:abordagens_pcg}

A literatura clássica de PCG classifica os métodos generativos em três grandes paradigmas operacionais: métodos construtivos, métodos baseados em busca (\textit{search-based}) e métodos baseados em restrições ou solucionadores (\textit{solver-based} / \textit{constraint-based}) \cite{pcgbook, summerville2018procedural, maleki2024procedural}.

Os métodos construtivos caracterizam-se por gerar o conteúdo em uma única passagem sequencial de passos algorítmicos, com tempo de execução previsível e sem realizar ciclos de avaliação e regeneração sobre o artefato produzido \cite{shaker2016constructive, summerville2018procedural}. Exemplos clássicos incluem funções de ruído coerente (como o ruído de Perlin para síntese de terrenos e texturas), autômatos celulares, algoritmos de particionamento espacial e gramáticas generativas \cite{linden2014dungeons, shaker2016constructive, maleki2024procedural}. Importa ressaltar que, embora procedimentos construtivos sejam computacionalmente rápidos e possam operar de forma determinística a partir de uma semente fixa, o determinismo algorítmico assegura apenas a repetibilidade da saída para uma mesma entrada; ele não garante, isoladamente, que o conteúdo resultante seja estruturalmente válido, conectado ou solucionável pelo jogador, a menos que as regras construtivas tenham sido estritamente projetadas para preservar tais invariantes em todas as etapas \cite{linden2014dungeons, shaker2016constructive}.

Os métodos baseados em busca (\textit{search-based PCG}) enquadram-se como um caso especial do paradigma gerar-e-testar (\textit{generate-and-test}) \cite{pcgbook, maleki2024procedural}. Em vez de aceitar imediatamente o artefato sintetizado ou simplesmente descartá-lo de forma binária, algoritmos baseados em busca atribuem a cada instância candidata um valor numérico contínuo por meio de uma função de avaliação ou aptidão (\textit{fitness function}), utilizando esse escore para guiar iterativamente a modificação e a seleção de novas soluções no espaço de busca \cite{linden2014dungeons, araujo2018geracao, maleki2024procedural}. Essa categoria emprega amplamente técnicas de computação evolutiva (como Algoritmos Genéticos) \cite{linden2014dungeons, maleki2024procedural}, além de meta-heurísticas como têmpera simulada (\textit{Simulated Annealing}), otimização por enxame de partículas (PSO) e Busca em Árvore de Monte Carlo (MCTS) \cite{maleki2024procedural, xu2025constraint}.

Os métodos baseados em restrições (\textit{constraint-based} ou \textit{solver-based PCG}), por sua vez, especificam declarativamente as propriedades e relações obrigatórias que o artefato final deve satisfazer, delegando a construção a solucionadores matemáticos ou lógicos (\textit{constraint solvers}), tais como programação por conjuntos de respostas (\textit{Answer Set Programming} --- ASP), solucionadores de satisfatibilidade módulo teorias (SMT) e bibliotecas de programação por restrições \cite{tanagra, linden2014dungeons, pcgbook, summerville2018procedural, chen2025narrative}, além de algoritmos de síntese baseados em propagação local de padrões e restrições, como o \textit{Wave Function Collapse} (WFC) \cite{maleki2024procedural}. Nessa abordagem, o foco desloca-se da definição de passos procedurais imperativos para a formulação rigorosa do espaço de design admissível \cite{tanagra, linden2014dungeons}.

\subsubsection{Geração Procedural de Fases}
\label{sec:pcg_fases}

A aplicação de PCG à criação de fases e níveis constitui um dos tópicos mais investigados na área, apresentando exigências técnicas particulares que a diferenciam da geração de ativos puramente cosméticos \cite{summerville2018procedural, maleki2024procedural}. A literatura distingue conteúdos cosméticos (como texturas e efeitos visuais, cuja alteração não modifica o resultado das ações do jogador) de conteúdos funcionais, que afetam diretamente a mecânica e a dinâmica da partida \cite{summerville2018procedural}. Fases de jogos são artefatos intrinsecamente funcionais: uma fase visualmente atraente, mas que apresente falhas topológicas ou bloqueios estruturais que impeçam sua conclusão, é considerada inválida para o propósito do jogo \cite{summerville2018procedural, linden2014dungeons}.

Na geração procedural de níveis e masmorras, um método típico envolve três componentes conceituais: (1) um modelo representacional abstrato que descreve a estrutura global da fase; (2) um método para construir esse modelo representacional; e (3) um método para sintetizar a geometria efetiva da fase a partir do modelo abstrato \cite{linden2014dungeons, shaker2016constructive}. Diversas famílias algorítmicas foram desenvolvidas para essas etapas. Algoritmos de particionamento de espaço, como o Particionamento Binário do Espaço (\textit{Binary Space Partitioning} --- BSP), \textit{quadtrees} (em 2D) e \textit{octrees} (em 3D), subdividem recursivamente o domínio espacial em células disjuntas, garantindo estruturalmente que salas alocadas em partições distintas não se sobreponham \cite{shaker2016constructive}. Em contrapartida, abordagens baseadas em agentes escavadores (\textit{agent-based dungeon growing}) operam em escala microscópica, deslocando um agente virtual que escava túneis e posiciona salas estocasticamente; embora produzam configurações orgânicas, agentes sem mecanismos de inspeção antecipada (\textit{look-ahead}) estão sujeitos a gerar salas sobrepostas e corredores sem saída \cite{shaker2016constructive}.

De modo complementar, autômatos celulares operam sobre grades discretas aplicando regras de transição baseadas na vizinhança local (como as vizinhanças de Moore ou de von Neumann) para formar cavernas naturais, mas não asseguram, por si sós, a conectividade global entre câmaras isoladas, exigindo etapas adicionais de verificação e abertura de passagens \cite{linden2014dungeons, shaker2016constructive}. Gramáticas generativas --- incluindo gramáticas de grafos e gramáticas de formas (\textit{shape grammars}) --- expandem um axioma inicial por meio de regras de produção e reescrita, permitindo capturar princípios de progressão e hierarquia em níveis de ação e aventura \cite{dormans2010, inproceedings, shaker2016constructive}. Apesar da maturidade dessas técnicas em domínios bidimensionais, a transição para a geração de fases tridimensionais complexas permanece um desafio aberto na literatura devido ao aumento expressivo das restrições geométricas, de navegação vertical e de acoplamento espacial \cite{linden2014dungeons, maleki2024procedural, xu2025constraint}.

% ---------------------------------------------------------------------------
\subsection{PCG Baseada em Aprendizado de Máquina e Uso de LLMs}
\label{sec:pcgml_llms}
% ---------------------------------------------------------------------------

\subsubsection{\textit{Procedural Content Generation via Machine Learning} --- PCGML}
\label{sec:pcgml}

A Geração Procedural de Conteúdo via Aprendizado de Máquina (\textit{Procedural Content Generation via Machine Learning} --- PCGML) é definida formalmente como a geração de conteúdos de jogos diretamente por modelos estatísticos ou de aprendizado de máquina que foram treinados sobre conjuntos de conteúdos de jogos preexistentes \cite{summerville2018procedural}. A distinção fundamental entre o PCGML e as abordagens clássicas baseadas em busca ou em solucionadores reside no mecanismo gerador: enquanto a busca evolucionária pode eventualmente empregar uma rede neural como função de avaliação para pontuar candidatos em um espaço de busca codificado manualmente, no PCGML o modelo treinado aprende a distribuição latente dos dados de treinamento e produz novas instâncias de conteúdo diretamente a partir de sua inferência \cite{summerville2018procedural}.

A taxonomia clássica de PCGML organiza os métodos segundo duas dimensões: a representação subjacente dos dados (sequências unidimensionais, grades bidimensionais/tridimensionais e grafos) e o paradigma de treinamento ou modelagem empregado \cite{summerville2018procedural}. Entre as principais técnicas investigadas na literatura figuram: (i) modelos probabilísticos baseados em contagem de frequência, como modelos de $n$-gramas, Cadeias de Markov Multidimensionais (MdMCs) e Campos Aleatórios de Markov (MRFs), utilizados para amostrar níveis fatia por fatia ou bloco por bloco; (ii) métodos de fatoração de matrizes e redução de dimensionalidade, como Análise de Componentes Principais (PCA) e fatoração de matrizes não negativas; (iii) Redes Neurais Recorrentes, em especial redes \textit{Long Short-Term Memory} (LSTMs), aplicadas à geração sequencial de fases de plataforma; (iv) \textit{Autoencoders} e \textit{Variational Autoencoders} (VAEs), capazes de comprimir níveis em espaços latentes contínuos para geração, interpolação e reparo de segmentos inválidos; e (v) Redes Adversárias Generativas (GANs) e Aprendizado por Reforço (PCGRL) \cite{summerville2018procedural, maleki2024procedural}.

Embora o PCGML reduza a necessidade de codificar manualmente regras de estilo ao extrair padrões diretamente de exemplos de design humano, a literatura aponta limitações estruturais importantes nesse paradigma \cite{summerville2018procedural, maleki2024procedural, her2025zeroshot3dmapgeneration}. A primeira refere-se à escassez crônica de dados (\textit{data scarcity}): diferentemente de domínios como visão computacional ou processamento de texto, jogos digitais dispõem de \textit{corpora} públicos reduzidos e carecem de formatos de representação padronizados entre diferentes títulos \cite{summerville2018procedural}. A segunda limitação diz respeito à dificuldade de assegurar restrições funcionais globais: modelos puramente estatísticos treinados sobre matrizes visuais ou sequências locais de blocos frequentemente reproduzem padrões estéticos plausíveis, mas falham em garantir propriedades de longo alcance, como a conectividade global e a solvabilidade do nível sem o auxílio de restrições externas \cite{summerville2018procedural, maleki2024procedural}.

\subsubsection{Modelos de Linguagem de Grande Porte Aplicados à Geração de Conteúdo}
\label{sec:llms_pcg}

Os Modelos de Linguagem de Grande Porte (\textit{Large Language Models} --- LLMs) constituem uma classe de modelos fundacionais de aprendizado profundo baseados na arquitetura \textit{Transformer}, introduzida para superar os gargalos computacionais e de memória dos modelos sequenciais precursores \cite{Gallotta_2024, maleki2024procedural}. Na evolução do Processamento de Linguagem Natural (PLN), abordagens estatísticas iniciais baseadas em cadeias de Markov e $n$-gramas deram lugar a representações vetoriais de palavras (\textit{word embeddings}) e a Redes Neurais Recorrentes (RNNs e LSTMs) \cite{Gallotta_2024, maleki2024procedural}. Contudo, como as arquiteturas recorrentes processam a informação sequencialmente passo a passo, elas sofrem de limitações na captura de dependências de longo alcance, associadas ao problema do desaparecimento de gradientes (\textit{vanishing gradient problem}) em sequências extensas \cite{maleki2024procedural}.

A arquitetura \textit{Transformer} contorna essas limitações ao substituir a recorrência sequencial pelo mecanismo de autoatenção (\textit{self-attention}) \cite{Gallotta_2024, maleki2024procedural}. Em vez de processar a entrada estritamente elemento por elemento, o mecanismo de autoatenção pondera dinamicamente a relevância mútua entre todos os elementos (\textit{tokens}) de uma sequência, permitindo ao modelo relacionar diferentes partes do contexto de entrada e capturar simultaneamente relações sintáticas, semânticas e estruturais de longo alcance \cite{Gallotta_2024, maleki2024procedural}.

Na literatura de modelos de linguagem baseados em \textit{Transformers}, destacam-se duas famílias principais: modelos bidirecionais de codificação (como o BERT, otimizados por modelagem de linguagem mascarada para extração de características e preenchimento contextual) \cite{Gallotta_2024, maleki2024procedural} e modelos autorregressivos (como a família GPT \cite{Gallotta_2024, maleki2024procedural} e modelos abertos como LLaMA e Mistral \cite{Gallotta_2024}, voltados à geração de texto e código). Na formulação autorregressiva causal, o modelo estima a distribuição de probabilidade condicional do próximo \textit{token} $x_t$ dado o contexto precedente $(x_1, x_2, \dots, x_{t-1})$ \cite{Gallotta_2024}. O desenvolvimento de um LLM moderno compreende tipicamente o pré-treinamento autossupervisionado em larga escala sobre vastos \textit{corpora} textuais, seguido pelo ajuste fino por instruções (\textit{instruction fine-tuning}) e pelo alinhamento com preferências humanas, notadamente via Aprendizado por Reforço com \textit{Feedback} Humano (RLHF) \cite{Gallotta_2024, her2025zeroshot3dmapgeneration}.

Em nível de inferência, o escalonamento para centenas de milhões ou bilhões de parâmetros faz emergir capacidades como o aprendizado em contexto (\textit{In-Context Learning} --- ICL) \cite{ocal2024sceneteller, her2025zeroshot3dmapgeneration}, pelo qual o modelo resolve novas tarefas condicionadas apenas pelas instruções do \textit{prompt} (\textit{zero-shot}) ou por poucos exemplos demonstrativos de entrada e saída (\textit{few-shot}), sem necessidade de atualização de seus pesos \cite{Gallotta_2024, ocal2024sceneteller, her2025zeroshot3dmapgeneration}. Estratégias complementares de engenharia de \textit{prompt} e acoplamento externo incluem a Cadeia de Pensamento (\textit{Chain-of-Thought} --- CoT), que decompõe problemas complexos em etapas explícitas de raciocínio intermediário, e a Geração Aumentada por Recuperação (\textit{Retrieval-Augmented Generation} --- RAG), que integra bases de conhecimento ou bancos de dados vetoriais externos ao contexto de geração \cite{Gallotta_2024, tang2025unrealllm, her2025zeroshot3dmapgeneration, wang2026designerlyloop}.

No domínio de PCG, a introdução dos LLMs mitigou parcialmente o problema da escassez de conjuntos de dados específicos de jogos ao transferir o conhecimento semântico e estrutural adquirido durante o pré-treinamento geral \cite{Gallotta_2024, todd2023level, maleki2024procedural}. A literatura documenta a aplicação de LLMs na geração de narrativas interativas, diálogos e árvores de missões; na síntese direta de níveis 2D codificados como grades de caracteres (a exemplo de \textit{Sokoban}, \textit{Super Mario Bros.} e mundos baseados em \textit{tiles}) \cite{todd2023level, nasir2024word2worldgeneratingstoriesworlds, maleki2024procedural}; e na atuação como planejadores semânticos de alto nível, capazes de traduzir intenções em linguagem natural em especificações de cena, parâmetros de ferramentas procedurais e regras de progressão \cite{Gallotta_2024, hu2024scenecraft, tang2025unrealllm, xu2025constraint, her2025zeroshot3dmapgeneration}.

\subsubsection{Geração de Representações Estruturadas}
\label{sec:representacoes_estruturadas}

Quando modelos de linguagem são integrados a sistemas computacionais mais amplos --- tais como motores gráficos, otimizadores matemáticos ou verificadores formais ---, a saída em linguagem natural livre torna-se inadequada devido à sua ambiguidade inerente e à ausência de garantias sintáticas para interpretação automatizada (\textit{parsing}) \cite{xu2025constraint, regmi2025system}. Para viabilizar a interoperabilidade entre o raciocínio probabilístico de um LLM e módulos algorítmicos determinísticos, a literatura recorre à geração de representações estruturadas intermediárias \cite{hu2024scenecraft, tang2025unrealllm, xu2025constraint, chen2025narrative, wang2026designerlyloop}.

Sob o ponto de vista teórico e de engenharia de software, é fundamental distinguir dois níveis conceituais nessa comunicação: o \textit{modelo conceitual} (ou estrutura abstrata de informação) e o \textit{formato de serialização} \cite{xu2025constraint, xu2025databasedrivenframework3dlevel, chen2025narrative}. O modelo conceitual corresponde à abstração matemática ou simbólica que organiza o domínio do problema, podendo assumir a forma de um grafo de cena, de raciocínio ou de progressão \cite{xu2025databasedrivenframework3dlevel, wang2026designerlyloop}, de um conjunto de triplas relacionais do tipo \texttt{[Objeto-Relação-Objeto]} \cite{chen2025narrative}, de uma árvore hierárquica de componentes, ou de sentenças formais em uma Linguagem de Descrição de Domínio Específico (DSL), como a \textit{Game Level Description Language} (GLDL) \cite{xu2025constraint} ou códigos de síntese programática \cite{hu2024scenecraft, yang2026codeasroomgenerating3drooms}.

O formato de serialização, por outro lado, consiste na notação textual padronizada utilizada para codificar, armazenar e transmitir esse modelo conceitual entre processos distintos \cite{sun2024roomplanner, xu2025databasedrivenframework3dlevel}. O formato JSON (\textit{JavaScript Object Notation}), assim como XML, YAML ou folhas de estilo estruturadas (CSS) \cite{ocal2024sceneteller}, constitui um exemplo amplamente adotado de representação serializada por estruturar dados em pares chave-valor e listas hierárquicas facilmente validáveis por esquemas formais (\textit{schemas}) \cite{tang2025unrealllm, sun2024roomplanner, xu2025databasedrivenframework3dlevel}. Portanto, o formato JSON não se confunde com o modelo lógico da fase em si: ele atua como o veículo sintático que codifica a abstração conceitual (por exemplo, os vértices, arestas e atributos de um grafo) para que sistemas externos possam processá-la sem ambiguidade \cite{tang2025unrealllm, xu2025databasedrivenframework3dlevel}.

% ---------------------------------------------------------------------------
\subsection{Raciocínio e Ancoragem Espacial em LLMs}
\label{sec:spatial_grounding}
% ---------------------------------------------------------------------------

\subsubsection{Raciocínio Espacial}
\label{sec:raciocinio_espacial}

O raciocínio espacial compreende a capacidade de representar mental ou computacionalmente objetos em um espaço dimensional, inferir suas propriedades geométricas e deduzir as consequências de suas relações e transformações mútuas \cite{sun2024roomplanner, xu2025constraint}. No contexto da síntese de ambientes virtuais e \textit{layouts} tridimensionais, o raciocínio espacial abrange tanto relações qualitativas e topológicas quanto relações métricas e contínuas \cite{xu2025constraint, chen2025narrative}.

As relações espaciais podem ser categorizadas em propriedades fundamentais que governam a organização de uma cena \cite{xu2025constraint, sun2024roomplanner, chen2025narrative}: (i) \textit{posição}, que define a localização exata de uma entidade em relação a um referencial de coordenadas global ou local; (ii) \textit{distância e proximidade}, que quantificam o afastamento métrico euclidiano entre pares de elementos (como estar próximo ou distante de um ponto de referência); (iii) \textit{direção e orientação}, que determinam para qual vetor ou ângulo a face frontal de um objeto está voltada (\textit{facing/focus}) e como seus eixos locais se alinham aos eixos do ambiente; (iv) \textit{adjacência e conectividade}, que indicam contato imediato ou comunicação entre regiões vizinhas (como duas salas que compartilham uma parede ou porta); (v) \textit{contenção}, que exige que objetos internos permaneçam estritamente dentro dos limites volumétricos de um compartimento maior; e (vi) \textit{interseção e ocupação espacial}, que descrevem se dois volumes tridimensionais compartilham pontos comuns no espaço \cite{xu2025constraint, sun2024roomplanner, xu2025databasedrivenframework3dlevel}. Manter a consistência simultânea de todas essas relações quando dezenas de elementos interagem em três dimensões constitui um problema combinatório e geométrico de elevada complexidade \cite{xu2025constraint, feng2026repurposing3dgenerativemodel}.

\subsubsection{\textit{Spatial Grounding}}
\label{sec:conceito_spatial_grounding}

Na interseção entre inteligência artificial simbólica, visão computacional e processamento de linguagem natural, o conceito de ancoragem (\textit{grounding}) refere-se ao estabelecimento de um vínculo semântico e operacional entre símbolos abstratos (como palavras, predicados ou \textit{tokens}) e seus referentes no mundo físico ou em um ambiente simulado \cite{Gallotta_2024, chen2025narrative}. Especificamente, a ancoragem espacial (\textit{spatial grounding}) consiste no processo de mapear descrições linguísticas e conceitos relacionais qualitativos --- tais como ``posicionar uma mesa de centro em frente ao sofá'', ``alocar uma estante encostada à parede norte'' ou ``conectar a sala de guarda ao corredor principal'' --- em parâmetros espaciais concretos, coordenadas métricas, escalas proporcionais e restrições físicas verificáveis dentro de uma representação geométrica 2D ou 3D \cite{hu2024scenecraft, sun2024roomplanner, chen2025narrative, ocal2024sceneteller}.

Uma ancoragem espacial efetiva exige que o sistema generativo não apenas compreenda a associação semântica e funcional (\textit{affordance}) entre os objetos mencionados no texto (por exemplo, saber que cadeiras costumam acompanhar mesas ou que portas conectam compartimentos) \cite{Gallotta_2024, chen2025narrative}, mas também traduza essas intenções em configurações que respeitem as dimensões reais dos modelos tridimensionais (\textit{bounding boxes}), as leis de suporte físico, a orientação funcional dos componentes e a preservação de áreas livres para circulação e interação \cite{sun2024roomplanner, xu2025constraint, ocal2024sceneteller}.

\subsubsection{Limitações dos LLMs em Tarefas Espaciais Tridimensionais}
\label{sec:limitacoes_espaciais_llms}

Em modelos generativos de linguagem, denomina-se ``alucinação'' a síntese de conteúdos que aparentam fluência e plausibilidade sintática, mas que são factualmente incorretos, contradictórios ou desprovidos de aderência às restrições do domínio \cite{Gallotta_2024, xu2025constraint, bandeali2026generative}. Quando LLMs são aplicados diretamente à geração de cenas e níveis tridimensionais --- sendo solicitados a produzir coordenadas numéricas $(x, y, z)$, dimensões e rotações para múltiplos ativos ---, esse fenômeno manifesta-se como \textit{alucinação espacial} (ou inconsistência geométrica) \cite{Gallotta_2024, sun2024roomplanner, yang2026codeasroomgenerating3drooms, feng2026repurposing3dgenerativemodel}.

A literatura recente analisa múltiplos fatores técnicos que explicam as limitações dos LLMs em tarefas espaciais 3D \cite{Gallotta_2024, hu2024scenecraft, sun2024roomplanner, feng2026repurposing3dgenerativemodel, yang2026codeasroomgenerating3drooms}. Embora os LLMs apresentem forte capacidade de planejamento semântico e extração de relações qualitativas, seu pré-treinamento ocorre sobre distribuições estatísticas textuais, sem acoplamento nativo a um verificador de geometria euclidiana ou motor de simulação física durante a inferência \cite{Gallotta_2024, sun2024roomplanner, feng2026repurposing3dgenerativemodel}. Além disso, ao tratar a síntese de \textit{layout} puramente como uma tarefa de geração textual, ocorre uma severa \textit{compressão de informação}: representar objetos tridimensionais complexos apenas por descritores textuais e caixas delimitadoras omite interações geométricas finas de superfície \cite{feng2026repurposing3dgenerativemodel}. Adicionalmente, em formulações autorregressivas ou em planejamentos de passo único (\textit{single-step planning}), verifica-se o fenômeno de \textit{viés de exposição} (\textit{exposure bias}) e acúmulo de erros: pequenas imprecisões métricas cometidas nas primeiras predições de coordenadas propagam-se ao longo da sequência gerada, degradando a coerência espacial global à medida que o número de objetos ou salas aumenta \cite{ocal2024sceneteller, sun2024roomplanner, feng2026repurposing3dgenerativemodel, yang2026codeasroomgenerating3drooms}.

Estudos empíricos na área de geração de cenas e níveis 3D documentam falhas recorrentes associadas a essa limitação de ancoragem espacial \cite{ocal2024sceneteller, hu2024scenecraft, sun2024roomplanner, xu2025constraint, feng2026repurposing3dgenerativemodel, yang2026codeasroomgenerating3drooms, bandeali2026generative}:
\begin{itemize}
    \item \textbf{Colisões e sobreposições volumétricas (\textit{overlapping} / \textit{clipping}):} atribuição de coordenadas cujos volumes delimitadores se interceptam ou se interpenetram, fazendo com que móveis, paredes ou salas inteiras ocupem simultaneamente a mesma região do espaço tridimensional \cite{ocal2024sceneteller, sun2024roomplanner, xu2025constraint, feng2026repurposing3dgenerativemodel};
    \item \textbf{Violações de suporte e de limites (\textit{out-of-bounds} e objetos flutuantes):} posicionamento de elementos fora do perímetro das salas, atravessando paredes estruturais ou suspensos no ar sem suporte físico adequado \cite{ocal2024sceneteller, sun2024roomplanner, xu2025constraint, feng2026repurposing3dgenerativemodel, bandeali2026generative};
    \item \textbf{Erros de orientação e alinhamento:} rotação incorreta de elementos funcionais (como portas, escadas ou mobiliário voltados contra paredes cegas) e desalinhamento entre aberturas de cômodos adjacentes \cite{sun2024roomplanner, xu2025constraint, yang2026codeasroomgenerating3drooms};
    \item \textbf{Obstrução topológica e perda de acessibilidade:} alocação de objetos ou estruturas que bloqueiam passagens críticas, portas ou corredores, fragmentando o espaço navegável e isolando itens essenciais à progressão \cite{sun2024roomplanner, xu2025constraint, xu2025databasedrivenframework3dlevel}.
\end{itemize}

Tais evidências na literatura demonstram que, embora os LLMs sejam ferramentas expressivas para a concepção semântica e estrutural de ambientes, a geração direta e irrestrita de coordenadas tridimensionais apenas por inferência textual mostra-se frequentemente insuficiente para assegurar a validade física e funcional de fases 3D, motivando o emprego de abstrações intermediárias, planejadores explícitos e algoritmos de otimização ou verificação espacial \cite{hu2024scenecraft, sun2024roomplanner, xu2025constraint, xu2025databasedrivenframework3dlevel, yang2026codeasroomgenerating3drooms}.

% ---------------------------------------------------------------------------
\subsection{Representação Lógica de Fases por Grafos}
\label{sec:grafos}
% ---------------------------------------------------------------------------

\subsubsection{Conceitos Fundamentais de Teoria dos Grafos}
\label{sec:conceitos_grafos}

A Teoria dos Grafos fornece o arcabouço matemático padrão para modelar relações discretas entre entidades \cite{summerville2018procedural}. Formalmente, um grafo é definido como um par ordenado $G = (V, E)$, no qual $V$ constitui um conjunto finito de vértices (ou nós) e $E$ representa um conjunto de arestas (ou arcos) que estabelecem ligações entre pares de vértices de $V$ \cite{maleki2024procedural, yang2026codeasroomgenerating3drooms}.

De acordo com a natureza de suas conexões, um grafo classifica-se como \textit{não direcionado}, quando suas arestas representam ligações simétricas e bidirecionais entre pares de vértices, ou como \textit{direcionado} (dígrafo), quando cada aresta possui uma orientação explícita que parte de um vértice de origem em direção a um vértice de destino \cite{dormans2010, summerville2018procedural}. O número de arestas conectadas a um vértice define o seu \textit{grau} \cite{inproceedings, xu2025databasedrivenframework3dlevel}. Adicionalmente, quando valores numéricos são associados às arestas, obtém-se um \textit{grafo ponderado}, no qual os pesos podem representar custos de travessia, distâncias ou níveis de dificuldade; quando símbolos, tipos ou conjuntos de propriedades são associados aos vértices ou arestas, tem-se um \textit{grafo rotulado} ou \textit{grafo de atributos} \cite{dormans2010, inproceedings, linden2014dungeons, chen2025narrative}.

Sobre essa estrutura definem-se conceitos topológicos essenciais para a análise de redes e níveis de jogos \cite{inproceedings, linden2014dungeons, summerville2018procedural}:
\begin{itemize}
    \item \textbf{Caminho (\textit{path}):} uma sequência finita de vértices conectados sucessivamente por arestas de $E$, permitindo transitar de um nó de origem a um nó de destino \cite{inproceedings, linden2014dungeons};
    \item \textbf{Ciclo (\textit{cycle}):} um caminho fechado que parte de um vértice e retorna a ele mesmo após percorrer uma sequência de arestas \cite{dormans2010, linden2014dungeons}. Um grafo desprovido de ciclos é denominado acíclico, destacando-se as \textit{árvores} (grafos conexos e acíclicos) e os Grafos Acíclicos Direcionados (DAGs) \cite{linden2014dungeons, wang2026designerlyloop};
    \item \textbf{Conectividade:} propriedade estrutural que indica a existência de caminhos ligando os vértices do grafo, assegurando que o conjunto de nós forme uma rede integrada e sem subcomponentes isolados \cite{linden2014dungeons, summerville2018procedural, xu2025databasedrivenframework3dlevel};
    \item \textbf{Alcançabilidade (\textit{reachability}):} propriedade que atesta se um vértice de destino pode ser efetivamente atingido a partir de um vértice inicial por meio de um caminho válido em $G$ \cite{dormans2010, linden2014dungeons, xu2025databasedrivenframework3dlevel};
    \item \textbf{Planaridade:} propriedade de um grafo que pode ser representado em um plano bidimensional de modo que suas arestas não se cruzem fora dos vértices onde incidem \cite{inproceedings}.
\end{itemize}

\subsubsection{Grafos Aplicados ao \textit{Level Design}}
\label{sec:grafos_level_design}

Na literatura de \textit{Level Design} e de geração procedural de masmorras, os grafos são amplamente adotados como modelos representacionais devido à sua capacidade de codificar relações estruturais de forma compacta \cite{dormans2010, inproceedings, linden2014dungeons, summerville2018procedural}. Essa modelagem manifesta-se principalmente sob duas perspectivas complementares: grafos topológicos de espaço e grafos de missão ou de ações \cite{dormans2010, inproceedings, linden2014dungeons}.

Nos \textit{grafos topológicos de espaço} (frequentemente chamados de grafos de salas ou redes de conectividade espacial), cada vértice $v \in V$ representa uma região física discreta do nível --- como uma sala, uma câmara, um pátio ou uma plataforma ---, enquanto cada aresta $e \in E$ representa uma ligação física de adjacência ou passagem entre duas regiões, tais como portas, corredores, pontes, escadas ou portais de teletransporte \cite{dormans2010, linden2014dungeons, summerville2018procedural}. Atributos associados aos vértices descrevem o tipo funcional ou temático da sala (ex.: sala inicial, sala de combate, sala de tesouro, sala de chefe) e os objetos nela contidos, ao passo que atributos nas arestas descrevem a natureza da transição (ex.: passagem aberta, porta trancada, passagem unidirecional) \cite{linden2014dungeons, xu2025databasedrivenframework3dlevel}.

Nos \textit{grafos de missão} ou \textit{grafos de ações} (\textit{mission graphs} / \textit{action graphs}), por sua vez, os vértices representam tarefas, eventos narrativos ou ações que o jogador deve realizar (por exemplo, ``derrotar inimigo'', ``coletar chave'', ``destrancar porta'', ``interagir com altar''), e as arestas direcionadas especificam as relações de precedência causal ou temporal entre essas ações \cite{dormans2010, inproceedings}. Trabalhos clássicos demonstram que um grafo de ações pode ser progressivamente transformado em um grafo de salas por meio do agrupamento de ações que devem ocorrer no mesmo local (\textit{co-located actions}), servindo como ponte formal entre o planejamento de \textit{gameplay} e a topologia do espaço \cite{inproceedings, linden2014dungeons}.

\subsubsection{Grafos de Progressão e Dependências}
\label{sec:grafos_progressao}

A combinação de topologia espacial com regras de precedência origina os \textit{grafos de progressão}, os quais modelam explicitamente os estados, pré-condições e dependências lógicas que governam o avanço do jogador através da fase \cite{dormans2010, inproceedings, linden2014dungeons, xu2025databasedrivenframework3dlevel}. O mecanismo canônico para expressar dependências em fases de aventura e exploração é a estrutura de chave e fechadura (\textit{lock-and-key puzzle}) \cite{dormans2010, linden2014dungeons, xu2025databasedrivenframework3dlevel}. Conceitualmente, uma ``fechadura'' representa qualquer barreira condicional associada a uma aresta ou nó (uma porta trancada, um abismo, um obstáculo destrutível ou um guardião), enquanto a ``chave'' corresponde ao recurso, item, habilidade ou estado (uma chave física, uma alavanca, um item como bombas ou um artefato especial) cuja obtenção prévia é pré-condição obrigatória para transpor a barreira \cite{dormans2010, inproceedings, shaker2016constructive}.

Em um grafo de progressão, essas dependências impõem restrições de ordem topológica ($\tau$): se uma aresta ou sala $R_L$ está condicionada a um bloqueio que exige um recurso situado na sala $R_K$, a precedência lógica requer que $R_K$ seja alcançável a partir do vértice inicial sem que seja necessário atravessar $R_L$ ($\tau(R_K) < \tau(R_L)$) \cite{xu2025databasedrivenframework3dlevel}. A organização dessas arestas e dependências define padrões estruturais importantes para o ritmo da fase: a rota crítica (caminho obrigatório entre o início e o objetivo final), as ramificações opcionais (subgrafos laterais que abrigam recompensas secundárias ou recursos auxiliares sem obstruir o fluxo principal) e os ciclos controlados, que permitem criar atalhos ou exigir retorno estratégico (\textit{backtracking}) após a obtenção de uma chave \cite{dormans2010, inproceedings, linden2014dungeons}.

\subsubsection{Estrutura Lógica versus Representação Espacial}
\label{sec:logica_vs_espacial}

Um aspecto teórico fundamental na modelagem por grafos reside no fato de que um grafo é uma estrutura puramente topológica e combinatória, e não uma representação métrica ou geométrica \cite{dormans2010, summerville2018procedural}. Os vértices e arestas de um grafo de progressão definem quais elementos existem, quais funções desempenham e como se conectam logicamente, mas não determinam, por si sós, coordenadas cartesianas $(x, y, z)$, dimensões em metros, formatos poligonais de salas ou comprimentos e curvaturas de corredores \cite{dormans2010, inproceedings, linden2014dungeons, summerville2018procedural}.

Essa abstração implica que a correspondência entre um grafo lógico e um \textit{layout} tridimensional é de um-para-muitos: um único grafo de progressão admite múltiplas realizações espaciais geometricamente distintas no espaço 3D, desde que as relações de adjacência e conectividade especificadas pelas arestas sejam preservadas \cite{dormans2010, inproceedings}. Ao mesmo tempo, nem todo grafo abstrato pode ser mapeado diretamente em uma grade ortogonal ou em um pavimento plano sem adaptações geométricas --- por exemplo, um grafo não planar ou com vértices cujo grau exceda o número físico de portas disponíveis em um módulo de sala requer a introdução de nós intermediários (como corredores ramificados ou múltiplos andares) durante a etapa de resolução de \textit{layout} \cite{inproceedings, linden2014dungeons, xu2025databasedrivenframework3dlevel}.

% ---------------------------------------------------------------------------
\subsection{Representação de \textit{Layouts} Espaciais 3D}
\label{sec:layout_3d}
% ---------------------------------------------------------------------------

\subsubsection{Representação Modular de Ambientes}
\label{sec:representacao_modular}

Para materializar um ambiente virtual a partir de uma especificação lógica, a literatura de PCG e a indústria de jogos empregam extensivamente a representação espacial modular \cite{linden2014dungeons, shaker2016constructive, xu2025databasedrivenframework3dlevel}. Nesse paradigma, o cenário tridimensional não é esculpido como uma malha única monolítica, mas construído pela composição de unidades geométricas discretas e reutilizáveis \cite{linden2014dungeons, shaker2016constructive, xu2025databasedrivenframework3dlevel}.

As unidades modulares variam em granularidade conforme a arquitetura do gerador \cite{linden2014dungeons, shaker2016constructive, vieira2024metodologia}:
\begin{itemize}
    \item \textbf{Grades e ladrilhos (\textit{grids}, \textit{tiles} e \textit{voxels}):} subdivisões regulares do espaço em células quadradas (2D) ou cúbicas (3D) de tamanho uniforme, nas quais cada célula recebe uma peça elementar (piso, parede, canto ou espaço vazio) \cite{pcgbook, shaker2016constructive, summerville2018procedural};
    \item \textbf{Blocos e segmentos (\textit{chunks}):} agrupamentos pré-autorados de elementos geométricos e interativos que formam seções parciais de um nível \cite{linden2014dungeons, shaker2016constructive};
    \item \textbf{Salas pré-construídas (\textit{room templates}) e \textit{prefabs} 3D:} módulos arquitetônicos completos (como salas temáticas de dimensões conhecidas, corredores e escadarias) ou ativos pré-fabricados (\textit{prefabs}) de objetos internos \cite{linden2014dungeons, xu2025databasedrivenframework3dlevel, regmi2025system}. Um \textit{prefab} encapsula em um único pacote reutilizável a malha poligonal tridimensional (\textit{mesh}), texturas, volumes de colisão física e metadados semânticos \cite{xu2025databasedrivenframework3dlevel, regmi2025system}.
\end{itemize}

Para viabilizar o acoplamento preciso entre módulos adjacentes (como unir uma sala a um corredor ou duas salas entre si), cada componente modular define pontos de conexão explícitos --- denominados conectores, âncoras, portas ou portais (\textit{connectors} / \textit{doorways} / \textit{exits}) \cite{inproceedings, linden2014dungeons, shaker2016constructive, sun2024roomplanner}. Esses conectores especificam a posição local da abertura na fronteira do módulo, suas dimensões de passagem e a direção ou orientação espacial para a qual a conexão se abre \cite{inproceedings, shaker2016constructive, sun2024roomplanner, xu2025databasedrivenframework3dlevel}.

\subsubsection{Relações e Transformações Espaciais}
\label{sec:transformacoes_espaciais}

A instanciação de componentes modulares em um espaço tridimensional contínuo fundamenta-se na geometria analítica e na álgebra linear, operando sobre dois referenciais distintos: o sistema de coordenadas local (intrínseco a cada módulo ou objeto, geralmente centrado em seu ponto de pivô ou origem geométrica) e o sistema de coordenadas global do mundo (\textit{world space}) \cite{sun2024roomplanner, xu2025constraint, ocal2024sceneteller}.

A configuração espacial (ou \textit{pose}) de cada componente $i$ no espaço global $\mathbb{R}^3$ é determinada por transformações afins que combinam três parâmetros fundamentais \cite{sun2024roomplanner, xu2025constraint}: (i) \textit{translação} (posição $\mathbf{p}_i = (x_i, y_i, z_i) \in \mathbb{R}^3$), que desloca a origem local do módulo para sua coordenada no mundo; (ii) \textit{rotação} (orientação $\mathbf{R}_i$, expressa por ângulos de Euler, matrizes de rotação ortogonais ou quatérnios), frequentemente parametrizada em \textit{layouts} arquitetônicos pelo ângulo de guinada $\theta_i$ em torno do eixo vertical para alinhar paredes e portas às direções cardeais \cite{sun2024roomplanner, xu2025databasedrivenframework3dlevel}; e (iii) \textit{escala} ou \textit{dimensão} ($\mathbf{s}_i = (w_i, l_i, h_i) \in \mathbb{R}^3$), que define a largura, o comprimento e a altura física ocupados pelo elemento \cite{sun2024roomplanner, xu2025constraint}.

Para simplificar o cálculo de ocupação espacial e de distâncias entre malhas poligonais arbitrárias, a computação gráfica e os sistemas de geração de \textit{layout} utilizam volumes delimitadores (\textit{bounding volumes}) tridimensionais \cite{sun2024roomplanner, xu2025constraint, xu2025databasedrivenframework3dlevel}. A representação padrão consiste na Caixa Delimitadora Alinhada aos Eixos (\textit{Axis-Aligned Bounding Box} --- AABB), definida pelos intervalos mínimos e máximos $[x_{\min}, x_{\max}] \times [y_{\min}, y_{\max}] \times [z_{\min}, z_{\max}]$ de cada objeto ao longo dos eixos coordenados, ajustada conforme a orientação espacial do elemento na cena \cite{sun2024roomplanner, xu2025constraint}.

\subsubsection{Restrições Geométricas}
\label{sec:restricoes_geometricas}

Para que um conjunto de módulos e objetos transformados no espaço $\mathbb{R}^3$ constitua um \textit{layout} arquitetonicamente íntegro, diversas restrições geométricas precisam ser satisfeitas simultaneamente \cite{linden2014dungeons, sun2024roomplanner, xu2025constraint, xu2025databasedrivenframework3dlevel}:
\begin{itemize}
    \item \textbf{Restrição de não sobreposição (\textit{non-overlap}):} dois módulos de salas, corredores ou objetos sólidos distintos não devem ocupar simultaneamente o mesmo volume no espaço contínuo; formalmente, a interseção volumétrica entre seus interiores (frequentemente quantificada pela métrica de Interseção sobre União tridimensional, $\text{IoU}_{3D}$, ou pelo volume de interseção $V_{\text{overlap}}$ de suas caixas delimitadoras) deve ser nula \cite{sun2024roomplanner, xu2025constraint};
    \item \textbf{Restrição de contenção e limites (\textit{boundary constraints}):} elementos internos alocados a uma sala devem ter seus volumes delimitadores integralmente contidos dentro do perímetro útil definido pelas paredes, piso e teto daquela sala, sem projetar-se para o exterior \cite{sun2024roomplanner, xu2025constraint};
    \item \textbf{Compatibilidade e alinhamento de conexões:} para que duas salas adjacentes ou uma sala e um corredor se conectem fisicamente, seus respectivos portais/portas devem coincidir em posição e dimensão, apresentando orientações compatíveis ao longo da fronteira compartilhada, sem que paredes sólidas obstruam a abertura \cite{inproceedings, shaker2016constructive, sun2024roomplanner, xu2025databasedrivenframework3dlevel};
    \item \textbf{Espaço livre e desobstrução (\textit{clearance}):} a disposição de objetos internos deve preservar corredores de circulação livres ao redor das entradas, saídas e pontos de interação, evitando o congestionamento excessivo (\textit{clustering}) que impeça a passagem física de agentes \cite{sun2024roomplanner, xu2025constraint, xu2025databasedrivenframework3dlevel}.
\end{itemize}

% ---------------------------------------------------------------------------
\subsection{Arquiteturas Híbridas e Geração Baseada em Restrições}
\label{sec:arquiteturas_hibridas}
% ---------------------------------------------------------------------------

\subsubsection{Sistemas Híbridos e Neuro-Simbólicos}
\label{sec:sistemas_hibridos}

A constatação de que nenhum paradigma isolado de PCG é isento de limitações --- métodos construtivos e baseados em restrições oferecem garantias estruturais, mas exigem codificação manual complexa de regras criativas, ao passo que modelos de aprendizado profundo e LLMs oferecem alta expressividade semântica, mas carecem de garantias espaciais e lógicas exatas --- impulsionou o desenvolvimento de arquiteturas híbridas e composicionais \cite{linden2014dungeons, Gallotta_2024, maleki2024procedural}. Na geração procedural composicional, métodos distintos são encadeados ou aninhados em um mesmo fluxo, de modo que as potencialidades de uma técnica compensem as fragilidades da outra \cite{linden2014dungeons, maleki2024procedural}.

Nas pesquisas contemporâneas que envolvem modelos fundacionais, essa integração assume frequentemente a forma de \textit{arquiteturas híbridas neuro-simbólicas} \cite{Gallotta_2024, hu2024scenecraft, chen2025narrative, xu2025constraint}. Nesses sistemas, o componente neural ou probabilístico (como um LLM) opera nas camadas de abstração semântica e simbólica --- interpretando intenções em linguagem natural, estruturando narrativas, definindo grafos de cena ou especificando restrições qualitativas ---, enquanto componentes simbólicos, formais ou algorítmicos (como solucionadores de restrições, algoritmos de busca em grafos, otimizadores matemáticos ou rotinas geométricas em motores gráficos) encarregam-se da resolução métrica, da instanciação espacial e da verificação rigorosa de regras \cite{hu2024scenecraft, tang2025unrealllm, chen2025narrative, xu2025constraint, xu2025databasedrivenframework3dlevel, her2025zeroshot3dmapgeneration}.

De maneira correlata, a literatura destaca os sistemas de \textit{iniciativa mista} (\textit{Mixed-Initiative PCG} ou cocriação humano-IA), exemplificados por ferramentas como \textit{Tanagra} \cite{tanagra} e \textit{Sentient Sketchbook} \cite{linden2014dungeons}, nos quais projetistas humanos e algoritmos generativos colaboram iterativamente sobre uma representação compartilhada \cite{tanagra, lai2020towards, walton2020evaluating, wang2026designerlyloop}. Nesses ambientes, o usuário pode especificar diretrizes de alto nível, curar cadeias de raciocínio intermediárias ou fixar partes da geometria, cabendo ao sistema procedural preencher, adaptar e validar o restante do nível mantendo o respeito ao controle do projetista e a fluidez do processo criativo \cite{tanagra, lai2020towards, bandeali2026generative, wang2026designerlyloop}.

\subsubsection{Geração Baseada em Regras e Restrições}
\label{sec:geracao_restricoes}

A geração controlada em sistemas formais e híbridos apoia-se teoricamente na modelagem por regras e na Satisfação de Restrições \cite{tanagra, linden2014dungeons, pcgbook, xu2025constraint}. Na formulação de um Problema de Satisfação de Restrições (\textit{Constraint Satisfaction Problem} --- CSP), o artefato a ser gerado é modelado por meio de um conjunto de variáveis de decisão (por exemplo, as coordenadas de posição, dimensões ou orientações atribuídas a cada componente do nível), de domínios numéricos ou discretos que delimitam os valores possíveis para cada variável (como as dimensões máximas da área de construção) e de um conjunto de restrições algébricas ou lógicas que especificam quais combinações simultâneas de valores são admissíveis \cite{tanagra, linden2014dungeons}.

Na formulação de problemas de \textit{layout} espacial e de design de fases, a literatura diferencia duas classes de restrições \cite{xu2025constraint, xu2025databasedrivenframework3dlevel}:
\begin{itemize}
    \item \textbf{Restrições rígidas (\textit{hard constraints}):} condições invioláveis cujo cumprimento integral é obrigatório para que a solução seja considerada admissível, tais como a ausência de sobreposição entre salas, a contenção de objetos dentro dos limites do mapa e a precedência topológica de uma chave em relação à porta que ela destranca \cite{tanagra, sun2024roomplanner, xu2025constraint};
    \item \textbf{Restrições flexíveis (\textit{soft constraints}):} preferências estilísticas, funcionais ou de qualidade espacial que não invalidam estruturalmente a fase caso sofram pequenos desvios, mas cujo atendimento deve ser maximizado, tais como distâncias ideais de proximidade entre objetos correlatos, grau de agrupamento (\textit{clustering}) ou distribuição equilibrada de espaços vazios (\textit{sparsity}) \cite{xu2025constraint, xu2025databasedrivenframework3dlevel}.
\end{itemize}

Quando o problema combina restrições rígidas e flexíveis no espaço contínuo, ele é frequentemente modelado como um Problema de Otimização de \textit{Layout} de Instalações (\textit{Facility Layout Problem} --- FLP), no qual uma função objetivo agrega termos de penalidade ponderados para cada violação de restrição, buscando uma configuração espacial que minimize o custo total \cite{xu2025constraint, xu2025databasedrivenframework3dlevel}.

\subsubsection{Validação e Reparo de Conteúdo Gerado}
\label{sec:validacao_reparo}

Como geradores estocásticos, heurísticos ou baseados em aprendizado de máquina podem produzir instâncias candidatas que violem uma ou mais restrições do domínio, a literatura de PCG estabelece diferentes estratégias teóricas para a garantia de qualidade e correção do conteúdo gerado \cite{pcgbook, summerville2018procedural, maleki2024procedural, xu2025databasedrivenframework3dlevel}:
\begin{enumerate}
    \item \textbf{Geração e teste com filtragem (\textit{generate-and-test} / rejeição):} o gerador produz um artefato completo ou parcial e o submete a rotinas determinísticas de verificação; caso alguma restrição rígida seja violada, a instância é integralmente rejeitada e uma nova geração é solicitada \cite{pcgbook, maleki2024procedural};
    \item \textbf{Busca construtiva com retrocesso (\textit{backtracking}):} durante a montagem passo a passo do nível (como na alocação sequencial de salas ou padrões geométricos), se o algoritmo detecta que a decisão atual inviabiliza a satisfação das restrições restantes, ele desfaz a última atribuição, poda o ramo inválido do espaço de busca e tenta uma configuração alternativa ou relaxa restrições locais \cite{tanagra, linden2014dungeons};
    \item \textbf{Reparo pós-geração (\textit{post-generation repair}):} em vez de descartar um nível quase válido, algoritmos especializados aplicam transformações locais corretivas sobre o artefato defeituoso, tais como o deslocamento incremental de objetos em colisão, a abertura algorítmica de passagens entre regiões desconectadas ou a remoção seletiva de elementos que obstruam portas e rotas de navegação \cite{summerville2018procedural, sun2024roomplanner, xu2025databasedrivenframework3dlevel};
    \item \textbf{Refinamento iterativo guiado por diagnóstico (\textit{feedback loop}):} quando o sistema conta com um planejador generativo ou módulo de otimização, os erros detectados pelos validadores formais (sintáticos, lógicos ou espaciais) são estruturados como sinais de diagnóstico para orientar uma nova iteração de ajuste da representação \cite{hu2024scenecraft, nasir2024word2worldgeneratingstoriesworlds, xu2025constraint, her2025zeroshot3dmapgeneration}.
\end{enumerate}

% ---------------------------------------------------------------------------
\subsection{Validade, Navegabilidade e Solvabilidade de Fases}
\label{sec:validade_navegabilidade_solvabilidade}
% ---------------------------------------------------------------------------

A avaliação formal de uma fase gerada proceduralmente exige decompor o conceito geral de ``qualidade funcional'' em propriedades distintas e verificáveis. Embora frequentemente relacionados na prática, os conceitos de validade lógica, validade geométrica, navegabilidade, busca de caminhos (\textit{pathfinding}) e solvabilidade operam em camadas diferentes de abstração e não são sinônimos \cite{dormans2010, inproceedings, linden2014dungeons, sun2024roomplanner, xu2025constraint}.

\subsubsection{Validade Lógica e Estrutural}
\label{sec:validade_logica}

A validade lógica e estrutural refere-se à consistência interna da fase em seu nível de representação abstrata ou topológica, independentemente de sua realização geométrica no espaço 3D \cite{dormans2010, inproceedings, linden2014dungeons}. Uma estrutura de nível é logicamente válida quando satisfaz as invariantes sintáticas e semânticas do domínio do jogo: a existência bem definida de um estado ou região inicial e de pelo menos uma condição ou região de objetivo final; a conectividade topológica do grafo subjacente (ausência de componentes obrigatórias isoladas); e a consistência causal das relações de dependência e pré-condições \cite{dormans2010, inproceedings, xu2025databasedrivenframework3dlevel}.

Violações de validade lógica ocorrem, por exemplo, quando o grafo de progressão contém dependências circulares insolúveis (\textit{deadlocks} lógicos) --- tais como posicionar uma chave exclusivamente no interior de uma sala cuja única porta de acesso requer aquela mesma chave para ser aberta ---, ou quando uma ramificação opcional absorve recursos indispensáveis sem permitir retorno ao caminho principal \cite{dormans2010, linden2014dungeons, xu2025databasedrivenframework3dlevel}.

\subsubsection{Validade Geométrica}
\label{sec:validade_geometrica}

A validade geométrica diz respeito à conformidade física e métrica da fase quando instanciada no espaço euclidiano bidimensional ou tridimensional \cite{sun2024roomplanner, xu2025constraint, feng2026repurposing3dgenerativemodel}. Um \textit{layout} 3D detém validade geométrica quando todas as salas, corredores e objetos internos estão posicionados dentro dos limites espaciais permitidos, sem interpenetrações ou sobreposições volumétricas inválidas (\textit{clipping}) entre malhas sólidas, e com alinhamento geométrico correto nas interfaces de conexão (como portais e batentes de portas entre módulos adjacentes) \cite{sun2024roomplanner, xu2025constraint, xu2025databasedrivenframework3dlevel}.

Na computação gráfica e em motores de jogos, a verificação da validade geométrica apoia-se em estruturas de dados e consultas espaciais, tais como testes de interseção entre caixas delimitadoras (como AABBs), árvores de particionamento espacial (\textit{quadtrees}, \textit{octrees} e árvores BSP) e lançamento de raios ou volumes (\textit{raycasting} e \textit{sphere casting}) \cite{shaker2016constructive, sun2024roomplanner, xu2025constraint, xu2025databasedrivenframework3dlevel}. Cabe destacar, sob o rigor conceitual, que estruturas como caixas delimitadoras, árvores espaciais e \textit{raycasting} não previnem nem corrigem, por si sós, sobreposições de malhas: elas constituem mecanismos matemáticos de consulta e detecção de colisão ou visibilidade, os quais fornecem os dados necessários para que algoritmos de posicionamento, otimização ou reparo atuem evitando ou eliminando configurações inválidas \cite{shaker2016constructive, xu2025constraint}.

\subsubsection{Navegabilidade}
\label{sec:navegabilidade}

A navegabilidade (\textit{navigability} ou \textit{traversability}) caracteriza a propriedade de um ambiente espacial de oferecer superfícies contínuas e fisicamente transponíveis para a locomoção de um agente virtual ou jogador, levando em conta suas restrições cinemáticas e dimensionais (como raio de colisão do corpo, altura vertical livre, inclinação máxima de rampas e altura máxima de degraus superáveis) \cite{sun2024roomplanner, xu2025constraint, xu2025databasedrivenframework3dlevel}. Mesmo que um cenário seja geometricamente válido (sem colisões entre paredes ou objetos), ele pode carecer de navegabilidade caso uma passagem seja estreita demais para o volume do agente, caso haja desníveis verticais sem escadas ou rampas compatíveis, ou caso móveis dispostos perto de uma porta obstruam o raio de passagem \cite{sun2024roomplanner, xu2025databasedrivenframework3dlevel}.

Para representar computacionalmente o subespaço navegável de um ambiente 3D, utilizam-se estruturas como grades de ocupação bidimensionais (\textit{occupancy grids}) \cite{sun2024roomplanner} e, de forma predominante em ambientes tridimensionais contínuos, as Malhas de Navegação (\textit{Navigation Meshes} --- \textit{NavMesh}) \cite{xu2025constraint, xu2025databasedrivenframework3dlevel, regmi2025system}. Uma \textit{NavMesh} consiste em uma representação geométrica das superfícies caminháveis (\textit{walkable surfaces} ou \textit{walkable areas}) do cenário tridimensional, delimitando as regiões livres de obstáculos estáticos sobre as quais agentes virtuais e jogadores podem se deslocar e executar algoritmos de busca de caminho com desvio de colisões \cite{xu2025constraint, xu2025databasedrivenframework3dlevel, regmi2025system}.

\subsubsection{\textit{Pathfinding}}
\label{sec:pathfinding}

A busca de caminhos (\textit{pathfinding}) consiste no processo algorítmico de encontrar uma rota contínua e livre de obstáculos entre uma posição de origem e uma posição de destino sobre uma representação discreta do espaço navegável, tal como uma grade de ocupação ou uma \textit{NavMesh} \cite{sun2024roomplanner, nasir2024word2worldgeneratingstoriesworlds, xu2025databasedrivenframework3dlevel}.

Enquanto algoritmos de busca em largura (\textit{Breadth-First Search} --- BFS) e de preenchimento por inundação (\textit{flood-fill}) exploram uniformemente os nós adjacentes para verificar a conectividade e identificar regiões conexas \cite{xu2025databasedrivenframework3dlevel}, a determinação de rotas em espaços métricos emprega amplamente algoritmos de busca informada por heurística, com destaque para o algoritmo $A^*$ (\textit{A-star}) \cite{sun2024roomplanner, nasir2024word2worldgeneratingstoriesworlds, xu2025databasedrivenframework3dlevel}. O algoritmo $A^*$ explora o espaço navegável combinando o custo acumulado desde o ponto de partida (como uma entrada ou porta) com uma estimativa heurística de distância até o destino para determinar se existe um trajeto viável e desobstruído entre dois pontos de interesse \cite{sun2024roomplanner, nasir2024word2worldgeneratingstoriesworlds, xu2025databasedrivenframework3dlevel}. No contexto de validação de níveis, a execução bem-sucedida de um algoritmo de \textit{pathfinding} sobre a grade de ocupação ou \textit{NavMesh} comprova a existência de conectividade física navegável entre as regiões e objetos do espaço \cite{sun2024roomplanner, xu2025databasedrivenframework3dlevel}.

\subsubsection{Solvabilidade de uma Fase}
\label{sec:solvabilidade}

A solvabilidade (\textit{solvability}, \textit{winnability} ou \textit{completability}) constitui o critério funcional mais abrangente de uma fase de jogo: ela atesta que existe ao menos uma sequência finita e executável de ações e deslocamentos que permite ao jogador transitar do estado inicial da fase até a condição de conclusão ou vitória, respeitando simultaneamente a geometria física do espaço e todas as regras, estados dinâmicos, pré-condições e restrições de mecânicas do jogo \cite{tanagra, linden2014dungeons, summerville2018procedural, nasir2024word2worldgeneratingstoriesworlds, xu2025constraint}.

A distinção teórica rigorosa entre validade geométrica, navegabilidade/\textit{pathfinding} e solvabilidade é indispensável porque o sucesso em uma dessas dimensões não implica automaticamente o sucesso nas demais \cite{dormans2010, linden2014dungeons, summerville2018procedural, xu2025constraint}:
\begin{itemize}
    \item \textbf{Validade geométrica sem navegabilidade:} duas salas podem estar posicionadas sem qualquer sobreposição de malhas (\textit{clipping} zero), mas separadas por uma falha na malha de piso ou por um obstáculo que impeça a formação de uma \textit{NavMesh} contínua entre elas \cite{sun2024roomplanner, xu2025databasedrivenframework3dlevel};
    \item \textbf{Navegabilidade e \textit{pathfinding} sem solvabilidade:} a existência de um caminho geométrico encontrado pelo algoritmo $A^*$ entre a sala inicial e a sala final sobre o piso estático não garante, por si só, que a fase seja solucionável. Se o caminho atravessar uma porta trancada cuja chave foi alocada após a própria porta, ou se os recursos exigidos para superar os desafios intermediários forem insuficientes, haverá conectividade espacial na malha geométrica, mas a fase será logicamente insolúvel no espaço de estados do jogo \cite{dormans2010, inproceedings, xu2025constraint};
    \item \textbf{Validade lógica abstrata sem solvabilidade espacial:} inversamente, um grafo de progressão pode ter todas as suas dependências de chaves e portas perfeitamente ordenadas no nível simbólico, mas resultar em uma fase insolúvel caso a etapa subsequente de instanciação geométrica 3D posicione um elemento que bloqueie fisicamente o corredor de acesso à sala da chave \cite{linden2014dungeons, xu2025databasedrivenframework3dlevel}.
\end{itemize}

Por conseguinte, a verificação plena da solvabilidade em fases com progressão condicional requer que a análise de alcançabilidade considere tanto o espaço de estados lógicos (transições condicionadas à posse de itens, chaves e cumprimento de objetivos) quanto a confirmação de que cada transição admitida no grafo lógico corresponde a uma rota fisicamente navegável no ambiente tridimensional gerado \cite{dormans2010, linden2014dungeons, xu2025constraint, xu2025databasedrivenframework3dlevel}.
